% !TeX root = ../数模通用模板.tex
\section{模型的评价、改进与推广}
\subsection{模型的优点}
首先，模型将预测、购电发布、储能执行和费用结算连接起来。不同发布时刻只使用当时可获取的信息，实际储电量逐槽递推并跨日连续，避免通过每日重置电池库存形成不真实的节费。供需误差不仅以单槽均方误差度量，还以整日路径进入调度，使连续缺口对储能的影响得到体现。

其次，费用与运行要求分层处理。问题一先确定满足运行设置的最低费用，再用明确的0.1\%费用让步改善功率变化与启动；问题二通过有限的计划模式变化约束组织充放电，再固定模式修正购电量。这些设置使费用代价与操作偏好可以分别检查，而非把任意惩罚直接计入购电账单。

最后，日内更新将新光伏信息和当前储电量共同用于剩余计划，并统一处理上调加价与下调退款。实际执行将购电盈余、可用库存和功率上限同时纳入充放电裁剪，剩余缺口才触发紧急购电。指定日期、年度汇总和工作簿均由同一执行轨迹产生，便于核对题目要求的结果。

\subsection{模型的局限}
预测与参数选择使用同一开发年度，当前结果说明在所给数据及信息条件下的表现，尚不足以保证跨年或不同小区的预测精度。价格采用点预测，其突变对购电与储能的联合影响未充分进入情景分布。

求解也存在范围限制。固定模式的购电修正采用局部优化；日内规划用分段短缺权重近似后续更新机会，未完整求解未来所有发布决策。各问的预测器、预报权限及运行设置并不完全相同，跨问费用差反映完整方案差异，不能全部归因于某一次预报或单独的价格变化。

平稳性指标与电池寿命之间仅有定性联系。反转次数、活动槽和等效循环能够描述运行强度，却未包含温度、荷电状态停留时间和循环深度等全部因素。部分动态工况对寿命的作用并非单调\cite{geslin2025}，因此本文不将减少某一指标直接换算为延寿或折旧收益。

\subsection{模型的改进}
后续可在独立年度或新小区上验证预测与调度参数，并通过保持预测器、初始状态和控制器相同的配对实验，分离价格变化与预报更新的作用。对波动电价，可进一步用已观察的价格误差构造联合情景；在通信与控制数据可得后，再评估测量延迟和执行偏差。若需要将寿命纳入费用，应引入可识别的循环老化模型及电芯参数\cite{xu2017aging}，据实际退化数据标定，不沿用当前操作惩罚作为折旧价格。

\subsection{模型的推广}
在保留负载保障、储电递推和发布时序的基础上，可将模型应用于园区光储系统、社区共享储能及充电站能量管理。具体应用时需重新确定外网容量、可转移负载、售电规则和价格可见范围；多储能单元可增加设备索引并分别施加功率、容量和效率约束。若系统存在孤岛运行或网络安全要求，还需增加备用、失负荷或潮流约束后重新求解。
